
\section{Discusión}

El resultado central de la Etapa~5 es una separación nítida entre lo que es pronosticable y lo que no lo es con la información disponible. El \textit{nivel} de afectación de cada cuadrante es estable y altamente explicable: los modelos de corte transversal alcanzan un $R^2$ de $0{.}90$ (SAIDI) y $0{.}84$ (SAIFI), significativo frente a una prueba de permutación, y conservan $0{.}42$ y $0{.}76$ cuando se predicen zonas no usadas en la estimación. La \textit{ocurrencia} de falla en la semana siguiente también es pronosticable con precisión útil (AUC de $0{.}84$). En cambio, la \textit{magnitud} de la afectación en una semana concreta, alrededor del nivel del cuadrante, no lo es: el mejor modelo de panel supera a la persistencia en un $20$--$26$~\%, pero no mejora de forma significativa el pronóstico ingenuo que consiste en usar el promedio histórico de cada cuadrante (Tabla~\ref{tab:diebold_mariano}), y su $R^2$ \emph{within} es inferior a $0{.}03$. En términos prácticos, la red de Norte de Santander tiene un componente territorial altamente predecible ---qué cuadrantes son sistemáticamente más vulnerables y cuáles tienen mayor probabilidad de falla la semana próxima--- y un componente de magnitud semanal que, con las variables disponibles, se comporta como ruido.

La comparación entre familias de modelos es coherente con esa lectura. Los efectos fijos de cuadrante son el mejor modelo de magnitud en ambos indicadores porque capturan exactamente el componente que es pronosticable, el nivel propio de cada cuadrante; por la misma razón, su desempeño es casi idéntico al de la media del cuadrante. Random Forest resulta inferior pese al ajuste de hiperparámetros, y la búsqueda selecciona la configuración más regularizada de la grilla, un patrón típico cuando la señal es débil y el modelo flexible tiende a ajustar ruido. Su sensibilidad a la cola de la distribución lo confirma: al recortar el $0{.}5$~\% de valores extremos del entrenamiento, su error se iguala al de los modelos lineales. Este resultado difiere del de los informes preliminares, en los que Random Forest obtenía el menor error en SAIDI. La diferencia se explica por dos errores de construcción de las variables sísmicas, detectados y corregidos en la re-estimación (Sección~\ref{sec:sismicidad_sgc}): con ellos, la última semana del periodo de prueba registraba un número de sismos veinticuatro veces superior al máximo de entrenamiento, y los modelos lineales extrapolaban ese valor de forma severa, mientras que el bosque, que no extrapola, quedaba protegido. Corregidos los datos, la aparente ventaja del modelo no lineal desaparece. El episodio ilustra por qué la verificación de cada variable contra su definición es parte del análisis y no un trámite previo.

Una lectura adecuada de estas cifras exige situarlas frente a la restricción de información bajo la cual fueron obtenidas. El $R^2$ de prueba de entre $0{.}20$ y $0{.}26$ en el panel semanal no es el máximo alcanzable con los datos disponibles, sino el máximo alcanzable \textit{sin incurrir en fuga de información}. La partición formal de las columnas del panel (Sección~\ref{sec:matrices_diseno}; Tabla~\ref{tab:particion_variables_modelado}) excluyó deliberadamente del conjunto de predictoras a \texttt{usuarios\_afectados} y \texttt{elementos\_afectados}, pese a que ambas presentan la asociación más alta del panel con los indicadores. La exclusión no obedece a un criterio estadístico sino de validez: \texttt{usuarios\_afectados} figura en el numerador de la definición de SAIDI y SAIFI, de modo que incorporarla como regresor equivaldría a explicar la respuesta mediante una transformación de sí misma, y produciría un ajuste próximo a la unidad sin ninguna capacidad de anticipación. La misma lógica excluyó los valores contemporáneos de \texttt{n\_eventos} y \texttt{duracion\_total} —admitidos solo en su versión rezagada— y la respuesta cruzada, que comparte el mismo numerador.

Esta decisión tiene una consecuencia que conviene enunciar sin ambigüedad: un modelo construido sin esa disciplina habría exhibido métricas notablemente superiores y habría sido, sin embargo, inútil en producción, porque en el momento de pronosticar la semana~$s$ el operador no dispone todavía de la información de severidad que solo se conoce al cerrar el evento. El desempeño reportado es, en ese sentido, una medida honesta de lo que el modelo puede ofrecer en las condiciones reales de uso, y constituye la referencia frente a la cual deberán compararse futuras iteraciones.

La autocorrelación espacial significativa detectada mediante el índice de Moran (SAIDI: $I=0{.}154$; SAIFI: $I=0{.}179$; ambos $p<0{.}001$) confirma que el fenómeno de calidad del servicio no es independiente entre cuadrantes vecinos, lo que valida la inclusión de rezagos espaciales y respalda, desde la evidencia empírica, el enfoque territorial adoptado desde el planteamiento del problema. Asimismo, las variables que resultan relevantes en al menos dos modelos (Tabla~\ref{tab:triangulacion}) ---el número de elementos y la longitud de red en ambos indicadores, y la vegetación y la proporción de clientes no residenciales en SAIDI--- coinciden con la literatura revisada en los antecedentes, que señala la infraestructura y la vegetación como factores recurrentes en la predicción de confiabilidad de redes de distribución. En este estudio, sin embargo, la señal de infraestructura resultó claramente más determinante que la climática, y la exposición sísmica, una vez corregida su construcción, solo conserva una asociación débil con SAIFI en el corte transversal.

La comparación con los antecedentes permite situar estos resultados. Primero, la débil señal del clima semanal contrasta con los estudios de huracanes y tifones, en los que la velocidad del viento es el predictor dominante~\cite{ref43,ref37}; la diferencia es coherente con la naturaleza de los datos, porque esos trabajos modelan eventos extremos de escala sinóptica, mientras que en Norte de Santander predominan tormentas convectivas locales de pocas horas que un acumulado semanal de reanálisis no captura, y porque la relación clima--interrupción documentada en la literatura es marcadamente no lineal~\cite{ref41}. Segundo, la elección de una celda de aproximadamente $5\times 2{.}5$~km se ubica en el rango de resoluciones que la literatura identifica como un compromiso entre señal y dispersión~\cite{ref68liu,ref41,ref42}, y la fracción de cuadrantes--semana sin eventos ($63$~\% en $D_{\mathrm{mod}}$) reproduce el problema de inflación de ceros que motivó los esquemas multietapa de~\textcite{ref42}, lo que respalda extender el modelo actual, condicional a la ocurrencia de falla, hacia una estructura de dos componentes. Tercero, la evidencia de que los modelos de árboles igualan o superan a especificaciones paramétricas con pocas variables~\cite{ref37} proviene de problemas con una señal meteorológica fuerte; cuando, como aquí, la señal dinámica es débil, la flexibilidad del modelo no se traduce en precisión, y la recomendación de~\textcite{ref43} de acompañar estos modelos con reentrenamiento periódico y con una referencia simple se vuelve especialmente pertinente. Finalmente, la heterogeneidad espacial encontrada coincide con la de~\textcite{ref8}, quienes muestran que los indicadores de confiabilidad varían en órdenes de magnitud entre zonas y periodos y que los valores agregados la ocultan: ese es precisamente el vacío de información que este trabajo aborda para el operador.

Desde el punto de vista de la gestión, la distinción entre variables relevantes y variables accionables condiciona el uso práctico de estos resultados. Las señales más robustas ---número de elementos y longitud de red--- pertenecen al ámbito de la planeación de inversión y son modificables en horizontes plurianuales; el NDVI es gestionable mediante programas de poda con ciclos de meses; la exposición sísmica no es intervenible y su utilidad se limita a la priorización territorial. Esta jerarquía, unida a los resultados de desempeño, precisa el aporte del trabajo al operador. No reside en pronosticar cuántas horas de interrupción tendrá un cuadrante la semana próxima, sino en dos productos de uso directo: el ordenamiento de los cuadrantes según su nivel estructural de afectación, para concentrar inversión, mantenimiento y poda, y la probabilidad semanal de falla por cuadrante, para programar cuadrillas y recursos de respuesta. Ambos son útiles para reducir los índices anuales y las compensaciones, que es el propósito del proyecto.

Entre las limitaciones del estudio, la más relevante es el sesgo de cobertura de emparejamiento evento--elemento detectado entre 2021 y 2023 (Sección~\ref{sec:resultados-etapa2}), que obligó a restringir el dataset de modelado a 2024--2025. Esta decisión, aunque necesaria para garantizar la validez de los resultados, redujo el horizonte temporal disponible para el entrenamiento y probablemente limita la capacidad de los modelos para capturar variabilidad interanual de más largo plazo; con solo $105$ semanas útiles, el \emph{holdout} del $15$~\% final deja una ventana de evaluación de aproximadamente dieciséis semanas, suficiente para detectar degradación sistemática pero no para contrastar la estabilidad del modelo ante un ciclo anual completo. Para mitigar esta limitación se complementó el \emph{holdout} con una validación de origen móvil en cuatro cortes (Tabla~\ref{tab:origen_movil}): el ordenamiento de los modelos se mantiene ---FE de cuadrante supera a la persistencia en los cuatro cortes y tiene el menor RMSE medio---, aunque la dispersión entre cortes es del mismo orden que la diferencia entre FE y Random Forest, lo que confirma que esa diferencia es pequeña.

En segundo lugar, la señal climática contemporánea resultó débil en el grano semanal (Sección~\ref{sec:resultados-etapa3}), lo que sugiere que futuras iteraciones del proyecto deberían profundizar en rezagos climáticos de mayor orden o en variables de exposición acumulada en lugar de solo variables climáticas instantáneas. Conviene precisar el alcance de esta limitación: lo que la evidencia descarta es la utilidad del clima \textit{simultáneo} como predictor en el grano semanal, no la existencia de un efecto climático sobre la continuidad del servicio. La agregación a semana completa promedia condiciones meteorológicas muy heterogéneas dentro del periodo y puede diluir el efecto de episodios intensos de pocas horas, que son precisamente los que la operación asocia a fallas; una verificación de esta hipótesis exigiría trabajar con el grano diario o con indicadores de intensidad máxima intrasemanal, no solo con acumulados.

En tercer lugar, el modelo de magnitud se estimó sobre las semanas con falla (\texttt{hubo\_falla}$=1$), decisión justificada por la cero-inflación del $74{.}9$~\% del panel (Sección~\ref{sec:segmentacion_analitica}). El componente de ocurrencia estimado en la re-estimación completa el esquema en dos partes, pero de forma básica: ambos componentes se estimaron por separado y su combinación no se calibró. Un modelo conjunto (por ejemplo, de tipo \emph{hurdle} o de conteo con inflación de ceros) permitiría evaluar el valor esperado incondicional con métricas propias de distribuciones con exceso de ceros.

En cuarto lugar, dos fuentes externas resultaron menos informativas de lo previsto. El registro de emergencias de la UNGRD disponible en datos abiertos termina en 2022 y no aporta información al periodo de modelado, y la sismicidad, una vez corregida, solo muestra una asociación débil en el corte transversal. Ambos hechos se documentan en el catálogo de datos como limitaciones de la fuente y no como evidencia de ausencia de efecto.

Finalmente, el modelo de efectos fijos de cuadrante y semana no resultó evaluable en el conjunto de prueba por falta de traslape entre semanas de entrenamiento y de prueba, una limitación conocida de los modelos con efectos fijos dobles en validaciones temporales estrictas (\emph{holdout}) y no un defecto del modelo en sí. Se reporta de forma explícita en las tablas de desempeño, antes que omitirse, porque documenta una incompatibilidad estructural entre efectos fijos temporales y validación prospectiva que cualquier replicación del trabajo volverá a encontrar.

\section{Trabajos futuros}
\label{sec:trabajo-futuro}

A partir de los hallazgos y las limitaciones identificadas, se proponen las siguientes líneas de trabajo futuro:

\begin{itemize}
    \item Resolver de forma retroactiva la incompatibilidad de esquemas de codificación entre BRAE y el inventario georreferenciado (Sección~\ref{sec:resultados-etapa2}), lo que permitiría extender el dataset de modelado más allá de 2024--2025 y evaluar la estabilidad de los modelos ante un horizonte temporal más largo.

    \item Explorar especificaciones espacio-temporales formales (por ejemplo, modelos de panel espacial dinámico o \emph{spatial Durbin}) que integren simultáneamente la dependencia espacial confirmada por el índice de Moran y la persistencia temporal identificada en el análisis exploratorio, en lugar de tratarlas como dos modelos separados (panel semanal y corte transversal).

    \item Incorporar variables de exposición climática acumulada (por ejemplo, precipitación acumulada a varias semanas o índices de sequía/humedad del suelo) dado que la señal climática contemporánea resultó débil, con el fin de evaluar si el efecto del clima sobre SAIDI/SAIFI opera con rezagos más largos que los considerados en este trabajo.

    \item Incorporar información de mayor resolución temporal que pueda explicar la magnitud semanal alrededor del nivel del cuadrante, que con los datos actuales se comporta como ruido: pronósticos meteorológicos de intensidad máxima intradiaria, alertas de tormenta eléctrica, programación de mantenimientos y maniobras, y disponibilidad de cuadrillas. Solo con información de este tipo tiene sentido volver a evaluar modelos no lineales más flexibles, como \emph{gradient boosting}.

    \item Consolidar el esquema en dos componentes en un modelo conjunto (\emph{hurdle} o de conteo con inflación de ceros), calibrar la probabilidad de falla y evaluar el valor esperado incondicional con métricas apropiadas para distribuciones con exceso de ceros, de modo que el resultado sea directamente utilizable para la programación operativa.

    \item Actualizar la fuente de emergencias UNGRD cuando se publiquen registros posteriores a 2022 e incorporar al flujo controles automáticos de cada variable integrada (rango, variación entre unidades y continuidad temporal), como los que permitieron detectar los errores de la variable sísmica.

    \item Contrastar la robustez de la partición de variables mediante un ejercicio de \emph{ablación} controlada: estimar deliberadamente una especificación que incluya el bloque de impacto posterior y cuantificar la magnitud exacta de la fuga de información. El resultado no sería un modelo utilizable, pero proporcionaría una cota superior del ajuste alcanzable y permitiría dimensionar con precisión el costo predictivo de la restricción \emph{ex ante} adoptada en este trabajo.

    \item Revisar los signos negativos estimados en el corte transversal para la longitud de red (\texttt{longitud\_red\_km}) y para la antigüedad del activo más viejo (Sección~\ref{sec:resultados-etapa5}), incorporando una medida explícita de densidad de usuarios por kilómetro de red y una fecha de instalación original de los activos, distinta de la última fecha de modificación registrada. Ambas hipótesis interpretativas son plausibles pero no verificables con la información actualmente disponible.

    \item Diseñar un esquema de actualización periódica del panel y de reentrenamiento de los modelos, de modo que el operador de red pueda adoptar el flujo documentado en este trabajo como una herramienta de monitoreo continuo y no solo como un ejercicio académico puntual, en línea con el objetivo de reproducibilidad planteado desde el inicio del proyecto.
\end{itemize}
